%! Intercepts, Zeros, and Extrema — Unit 2, Lesson 2.4 — Exit Ticket (Answer Key)
\documentclass[10pt]{article}
\usepackage{atda-article}
\usepackage{atda-key}   % pulls in -boxes; do NOT also load -boxes
\newcommand{\TallMath}[1]{$\displaystyle #1\rule[-1.4em]{0pt}{3.2em}$}
\setlength{\workrowsep}{8pt}

\begin{document}

\pageheader{Unit 2, Lesson 2.4}{Exit Ticket --- Answer Key}
% NAMESTRIP: no \namedateperiod here — mirrors the blank. Teacher-only prose goes
% in the lesson plan as \begin{teachernote}[Exit Ticket], never in a key.

\vspace{0.15in}

\begin{notesbox}{One Creek, and One Claim in the Wrong Units}
\small
Notes closed. Three items. The graph shows $L(t)$, the water level of Miller's Creek in inches
\emph{above or below its normal level}, $t$ days after Monday.

\begin{center}
\begin{tikzpicture}
\begin{axis}[
    width=10.2cm, height=4.6cm,
    xlabel={$t$ (days after Monday)}, ylabel={$L$ (inches from normal)},
    xmin=0, xmax=7, ymin=-4.4, ymax=10.4,
    xtick={0,1,2,3,4,5,6,7}, ytick={-4,-2,0,2,4,6,8,10}, minor y tick num=1,
    grid=both, grid style={linegray!45}, minor grid style={linegray!30},
    axis lines=left,
    tick label style={font=\tiny}, label style={font=\scriptsize}]
  \addplot[royal, very thick] coordinates {(0,6) (3,-3) (7,9)};
  \addplot[keyred, only marks, mark=*, mark size=1.7pt]
    coordinates {(0,6) (2,0) (4,0) (3,-3) (7,9)};
\end{axis}
\end{tikzpicture}
\end{center}

\begin{enumerate}[leftmargin=1.6em, itemsep=7pt, topsep=4pt]

\item Give the $y$-intercept as a point, give the two zeros of $L$ as numbers, and say what a zero
means about this creek. Then name the interval where $L$ is decreasing and the interval where it is
increasing.

\ansline{$y$-intercept $(0,6)$. \quad Zeros: $t=2$ and $t=4$ --- both are \emph{days}, not levels.}\\
\ansline{A zero means the creek is at exactly its normal level that day, $0$ inches off.}\\
\ansline{Decreasing on $(0,3)$; increasing on $(3,7)$. Both intervals are days.}\\

\item Give the absolute maximum and the absolute minimum of $L$ over the week, each as a
\emph{value and a location}. For each one, say whether it happens at a turning point or at an
endpoint.

\ansline{Absolute maximum $9$ inches, at $t=7$ --- at an \textbf{endpoint}, not at a turn.}\\
\ansline{Absolute minimum $-3$ inches, at $t=3$ --- at the \textbf{turning point}.}\\

\item A classmate writes: \emph{``The water level is increasing from $-3$ to $9$.''} Both of those
numbers are answers you just wrote down in item 2. Explain what is wrong with the sentence anyway,
and write the correct statement.

\ansline{$-3$ and $9$ are water \emph{levels} --- outputs, read on the vertical axis. They are the}\\
\ansline{right numbers for item 2 and the wrong ones here: ``increasing'' asks \emph{which days}.}\\
\ansline{Correct: $L$ is increasing on $(3,7)$, and over those days it rises from $-3$ to $9$ inches.}\\

\end{enumerate}
\end{notesbox}

\end{document}
